\chapter{Literature Review}
\label{ch:literature}
\newpage

\section{Foundations of Software Quality and Design}
\label{sec:foundations}

Evaluating coding agents against design criteria requires a clear account of those criteria. This
section covers design principles, code smells and the metrics used to detect them.

\subsection{Design Principles}
\label{sec:design-principles}

Robert C. Martin's design principles \citep{martin2000solid,martin2003agile}, grouped under the
acronym SOLID, are the most widely taught framework for object-oriented design.
As Figure~\ref{fig:solid-wheel} shows, each principle can be pictured through the typical shape of
its violation (in the interface card, dotted members are stubs). The Single Responsibility Principle (SRP) states that a class should
have one reason to change. The Open/Closed Principle (OCP) states that a module should be open for
extension and closed for modification. The Liskov Substitution Principle (LSP) requires that a
subtype can be used wherever its base type is expected. The Interface Segregation Principle (ISP)
states that clients should not depend on methods they do not use. The Dependency Inversion
Principle (DIP) states that high-level modules should depend on abstractions, not on concrete
details. These are heuristics, but they connect to a long line of metrics work: cohesion, coupling
and inheritance depth can all be measured, which is what makes automated assessment possible.
Larman's GRASP patterns \citep{larman2004grasp} give finer guidance on assigning responsibilities,
such as Information Expert, Low Coupling and High Cohesion.

\begin{figure}[htb]
\centering
% Figure: the five SOLID principles around a central hub.
%
% Layout note: the five boxes sit on a TRUE pentagon, centres at radius 3.25 and
% angles 90, 162, 234, 306, 18 degrees, so the figure is symmetric about the
% vertical axis. The earlier version placed them by eye (top box at y=3.3 but
% the bottom pair at y=-2.9), which read as lopsided. Every box carries the same
% text width and minimum height, and each tagline is short enough to sit on one
% line, so no box grows taller than its neighbours.
\begin{tikzpicture}[
  every node/.style={font=\footnotesize},
  box/.style={draw,thick,rounded corners=3pt,align=center,text width=3.15cm,
              minimum height=1.55cm,font=\scriptsize,inner sep=4pt},
  spoke/.style={-{Latex[length=2mm]},thick,shorten >=2pt},
]
\node[circle,draw=issatblue,very thick,fill=lightblue,minimum size=1.9cm,
      font=\bfseries,text=issatblue] (hub) at (0,0) {SOLID};

% pentagon: r = 3.25, starting at 90 degrees and stepping by 72
\node[box,draw=srpcolor,fill=srpcolor!10] (srp) at (90:3.25)
  {\textbf{\textcolor{srpcolor}{S: SRP}}\\Single Responsibility\\[2pt]\textit{one reason to change}};
\node[box,draw=dipcolor,fill=dipcolor!10] (dip) at (162:3.25)
  {\textbf{\textcolor{dipcolor}{D: DIP}}\\Dependency Inversion\\[2pt]\textit{depend on abstractions}};
\node[box,draw=ispcolor,fill=ispcolor!10] (isp) at (234:3.25)
  {\textbf{\textcolor{ispcolor}{I: ISP}}\\Interface Segregation\\[2pt]\textit{no unused members}};
\node[box,draw=lspcolor,fill=lspcolor!10] (lsp) at (306:3.25)
  {\textbf{\textcolor{lspcolor}{L: LSP}}\\Liskov Substitution\\[2pt]\textit{subtypes substitutable}};
\node[box,draw=ocpcolor,fill=ocpcolor!10] (ocp) at (18:3.25)
  {\textbf{\textcolor{ocpcolor}{O: OCP}}\\Open/Closed\\[2pt]\textit{extend, don't modify}};

\draw[spoke,srpcolor] (hub) -- (srp);
\draw[spoke,dipcolor] (hub) -- (dip);
\draw[spoke,ispcolor] (hub) -- (isp);
\draw[spoke,lspcolor] (hub) -- (lsp);
\draw[spoke,ocpcolor] (hub) -- (ocp);
\end{tikzpicture}

\caption{The five SOLID principles.}
\label{fig:solid-wheel}
\end{figure}

Evidence on how well agents follow these principles is recent. The NITR benchmark
\citep{zhu2026nitr} tests whether agent edits keep code maintainable, over 21 curated C\texttt{++}
repository probes spanning nine maintainability dimensions, and finds dependency control (4.3\%
success) harder than responsibility decomposition (15.2\%). \citet{pehlivan2025solid}
prompt four language models to detect SOLID violations in 240 manually validated examples: Single
Responsibility is the easiest principle to detect, Dependency Inversion the hardest, and accuracy
drops as code becomes more complex.

Much of the agent-evaluation work reviewed in this chapter, including most of the 2026 studies in
Table~\ref{tab:comparative-synthesis}, is available only as preprints. These studies are cited
because they are the closest related work. Where the thesis relies on one, it names the specific
result it uses.

\subsection{Code Smells and Their Detection}
\label{sec:code-smells}

\citet{fowler1999refactoring} popularised the term \emph{code smell} for patterns such as Long
Method, God Class, Data Clumps and Feature Envy, which point to design problems without being bugs.
\citet{suryanarayana2014smells} later catalogued design smells by the design principle they violate;
for example, \emph{Missing Hierarchy} is conditional logic on type codes used where a class
hierarchy could hold the variation. DesigniteJava is a tool that detects smells from that
catalogue.

The link between smells and maintenance cost has been studied but is not uniform.
\citet{khomh2012antipatterns} find that classes participating in antipatterns are more change-prone
and fault-prone than other classes. \citet{olbrich2010godclass} find that God Classes and Brain
Classes change more often and contain more defects, but that after normalising for size they change
less often and contain fewer defects than other classes. Size alone is therefore a weak indicator of
a design problem, and a detector that relies on it should be read with caution.

Automated detection usually works through metric thresholds. \citet{marinescu2004detection}
formalised this as \emph{detection strategies}: rules that combine several metrics, each compared
with a threshold, to locate a specific design flaw such as a God Class. Tools such as PMD,
SonarQube, CodeScene and the CK metrics suite apply rules of this kind to many codebases at once.

\citet{zhu2026smells} study smells in AI-generated code. They report a \emph{Volume-Quality Inverse
Law}, in which code volume is a near-perfect predictor of structural degradation, and a
\emph{Reasoning-Complexity Trade-off}, in which more capable models generate more bloated and more
coupled code. They conclude that addressing this requires explicit architectural foresight during
generation.

\subsection{Metrics}
\label{sec:architectural-metrics}

The Chidamber and Kemerer (CK) suite \citep{chidamber1994metrics} defines Weighted Methods per Class
(WMC), Depth of Inheritance Tree (DIT), Number of Children (NOC), Coupling Between Objects (CBO),
Response For a Class (RFC) and Lack of Cohesion in Methods (LCOM). LCOM has several later variants,
including the connected-components form LCOM4 \citep{hitz1995cohesion} and the normalised LCOM*
\citep{hendersonsellers1996metrics}. For packages, Martin's instability and abstractness metrics
\citep{martin2003agile} measure how much a package depends on others and how abstract it is.

RealBench \citep{li2026realbench} measures architectural quality in repository-level generation.
Its 61 repositories come with natural-language requirements and UML designs, and generated code is
scored on an Architecture@$k$ metric from 0 to 4. The best score across the models studied is
2.29. Test pass rates fall with repository size, from 25.48\% for repositories of up to 500 lines to
5.81\% for repositories of 2,000 lines or more.

\section{Evolution of Coding Agents}
\label{sec:evolution}

AI coding tools have changed quickly (Figure~\ref{fig:timeline}, p.~\pageref{fig:timeline}). The
newer tools take on different kinds of task, and this affects how they should be evaluated.

\subsection{First-Generation Assistants (2021--2022)}

GitHub Copilot, released in 2021 and built on OpenAI's Codex model, suggested the next line or
function from the surrounding code. Amazon CodeWhisperer (2022) added security scanning and
reference tracking for suggestions that resemble licensed code. Both reacted to what the developer
was typing. HumanEval \citep{chen2021codex}, which measures pass@$k$ on 164 hand-written Python
problems, suited this kind of tool.

\subsection{Agentic Tools (2023--2025)}

Devin (Cognition, 2024) was presented as an agent that completes tasks end to end, using a shell, a
browser and an editor. Cursor added repository indexing and agentic editing to an IDE. Claude Code
(Anthropic, 2025) runs as an agent in the terminal. OpenCode (2025) is an open-source agent, and
OpenAI's Codex CLI (2025) runs a local agent loop with shell access.

These tools keep state across a session, break a task into steps, call tools and retry after
failures. A function-level benchmark such as HumanEval does not exercise these abilities.

Few peer-reviewed comparisons of these tools exist. Practitioner reports describe design choices
inside them. Anthropic's report on context engineering \citep{anthropic2025context} describes a
language model's context as a limited attention budget and names bloated tool sets, which leave the
model unsure which tool to use, as one of the most common failure modes it observes. The report
comes from the vendor of one of the agents measured in this thesis, so it is cited only for this
design argument and not as evidence. \citet{bensghaier2026harness} provide empirical support: with
the model held fixed, quality changes across releases of the agent software.

\subsection{Functional Benchmark Performance}

On function-level benchmarks agents score highly. AgentCoder \citep{huang2023agentcoder}, which
uses separate programmer, test designer and test executor agents, reaches 96.3\% pass@1 on
HumanEval with GPT-4. HumanEval tests single functions, where design problems are least visible.

Harder benchmarks give lower scores. BigCodeBench \citep{zhuo2024bigcodebench} contains 1,140 tasks
that require calls to 139 libraries. Across about 60 models the best score is 60\%, against 97\% for
humans. LiveCodeBench \citep{jain2024livecodebench} uses problems published after a model's
training cut-off in order to detect contamination. SWE-EVO \citep{sweevo2025} contains 48
long-horizon evolution tasks that change 21 files on average. GPT-5.4 with OpenHands resolves 25\%
of them, while GPT-5.2 resolves 72.80\% of SWE-bench Verified.

\section{Evaluation of Design Quality in AI-Generated Code}
\label{sec:evaluation-review}

A growing body of work examines structural properties of agent output. This section groups it by
what it measures.

\subsection{Principle Adherence}

\citet{pehlivan2025solid} evaluate language models as detectors of SOLID violations (see
Section~\ref{sec:design-principles}). \citet{beyond2026designissues} generate large projects with
Cursor, averaging 16,965 lines of code and 114 files. The projects reach 91\% functional correctness
on average, yet CodeScene reports 1,305 design issues and SonarQube 3,193, most often code
duplication, high complexity and large methods, which the authors relate to SRP and to separation
of concerns.

\citet{konrad2026vibe} describe how agents make architectural decisions. Agents such as Claude
Code and Codex choose frameworks, configure databases and set up authentication, and the authors
treat each of these choices as an architectural decision. Without explicit rules, agents fall back
on patterns learned from training data, and the decisions leave no design documents behind. This
suggests that agent-written code can end up with an unplanned structure.

\subsection{Code Smells in Agent Output}

\citet{zhu2026smells} argue that AI-generated code carries a distinct ``machine signature'' of
defects, including broken abstractions (an interface created and then bypassed through concrete
types), systematic duplication, over-engineering and verbosity. Broken abstractions relate to the
Dependency Inversion and Open/Closed principles.

Some design problems have a structural form and can be detected in principle. Others concern intent:
code with an unremarkable shape that was not written to any coherent design. A static rule cannot
detect the second kind, so a detector based on structure has to be checked against a reader of the
code before its flags are treated as verdicts on design.

\subsection{Repository-Scale and Acceptance Studies}

ProjDevBench \citep{lu2026projdevbench} evaluates six coding agents on 20 project-development
problems in eight categories, combining online-judge testing with LLM-assisted code review. The
overall acceptance rate is 27.38\%. Agents handle basic functionality and data structures but
struggle with complex system design, time complexity and resource management. With 20 problems, a
comparison between individual agents is limited.

Two studies use the AIDev corpus that this thesis also uses. \citet{mazloomzadeh2026agentic}
compare agentic and human pull requests over successive quarters: how their merge rates differ, which
tasks agents are given and how the characteristics of the requests change. \citet{cynthia2026postmerge}
analyse 1,210 merged agent bug-fix pull requests from Python repositories with SonarQube, comparing
the base and merged commits so that only issues introduced by the pull request are counted. Three of
their results bear on this thesis. Differences between agents in raw issue counts largely disappear
after normalising by code churn. Code smells are more common than bugs, especially at critical and
major severity. And merge success does not reliably reflect post-merge code quality.

\citet{sawada2026maintenance} study what happens to agent-generated code after it is merged, over
more than 1,000 files and about 3,200 changes in 100 popular repositories from AIDev. AI-generated
files are maintained less often than human-written files, and each update touches a small fraction
of the file. Updates to AI-written files are mostly feature extensions, updates to human-written
files mostly bug fixes, and human developers carry out most of the maintenance in both cases. Agent
code is therefore extended by people other than its author, which is the situation in which its
structure matters.

SWE-bench \citep{jimenez2024swebench} measures how many GitHub issues a system resolves with the
associated tests passing. It evaluates repair of single issues and does not evaluate design.
\citet{metr2026sweverified} show how far test-based scores are from maintainer judgement: roughly
half of the SWE-bench Verified patches that pass the tests would not be merged, and maintainer
decisions are about 24 percentage points below the automated score. Because the agents could not
respond to review feedback, the authors do not present the gap as a fixed capability limit.

SlopCodeBench \citep{slopcodebench2026} follows 15 agents through 36 problems and 196 checkpoints
in which each task extends the agent's own earlier solution. No agent solves a problem end to end,
and the best agent passes 14.8\% of checkpoints. Structural erosion increases in 77\% of
trajectories and verbosity in 75.5\%, and agent code is 2.3 times more verbose than comparable
open-source Python code.

\subsection{Prompting and Mitigation}

Prompt design changes detection quality: in \citet{pehlivan2025solid} an ensemble prompt works best
for OCP and an example-based prompt for DIP. For generation, SlopCodeBench reports that
quality-aware prompts reduce initial verbosity and erosion by up to a third but do not change the
rate at which quality degrades \citep{slopcodebench2026}.

Linting loops, in which an agent runs a checker and revises its code before committing, are
recommended in practitioner writing. We found no controlled study of their effect on design
quality, as opposed to lint compliance. Such a loop can only enforce what its checker can express, and a checker based on structure cannot
express design intent.

\section{Comparative Synthesis}
\label{sec:comparative-synthesis}

As Table~\ref{tab:comparative-synthesis} shows, the most closely related studies differ in what they
evaluate and measure, and each stops short of assessing design quality across several agents on real
contributions. The last column of the table states where each study's scope stops.

\begin{landscape}
\footnotesize
\begin{longtable}{@{}L{3.0cm}L{2.6cm}L{3.4cm}L{4.7cm}L{4.7cm}@{}}
\caption{Closely related studies.}
\label{tab:comparative-synthesis}\\
\toprule
\textbf{Study} & \textbf{Systems evaluated} & \textbf{Measures} & \textbf{Main findings} & \textbf{Limitation} \\
\midrule
\endfirsthead
\multicolumn{5}{l}{\footnotesize\itshape (Table~\ref{tab:comparative-synthesis} continued)}\\
\toprule
\textbf{Study} & \textbf{Systems evaluated} & \textbf{Measures} & \textbf{Main findings} & \textbf{Limitation} \\
\midrule
\endhead
\bottomrule
\endfoot
\bottomrule
\endlastfoot
\citet{pehlivan2025solid} & 4 LLMs & LLM detection of SOLID violations (240 examples) & SRP easiest to detect, DIP hardest; accuracy falls with complexity & Detection on examples, not measurement of agent contributions \\
\addlinespace
\citet{beyond2026designissues} & Cursor & Design issues in generated projects & 91\% functional correctness, yet 1,305 CodeScene and 3,193 SonarQube issues & One tool; no comparison between agents \\
\addlinespace
\citet{konrad2026vibe} & Several agents & Implicit architectural decisions & ``Vibe architecting''; one task gave 141 to 827 lines over 2 to 6 files & Qualitative; no measurement over a corpus \\
\addlinespace
\citet{lu2026projdevbench} & 6 coding agents & Online-judge tests and LLM code review & 27.38\% acceptance; weak on system design & 20 problems; no per-principle measurement \\
\addlinespace
\citet{jimenez2024swebench} & Language models and agents & Issues resolved with tests passing & Standard issue-level benchmark & Issue repair only; no design measurement \\
\addlinespace
\citet{metr2026sweverified} & Agents on SWE-bench Verified & Maintainer merge decisions & About half of passing patches would not be merged; 24-point gap & Merge verdict only \\
\addlinespace
\citet{zhu2026nitr} & Agents & Maintainability tasks with a structural oracle, 21 C\texttt{++} probes & 4.3\% dependency control, 15.2\% responsibility decomposition; 13.3\% pass tests but fail the oracle & Benchmark tasks in one language, not contributions to real projects \\
\addlinespace
\citet{slopcodebench2026} & 15 agents & Verbosity and erosion across checkpoints & Best agent passes 14.8\% of checkpoints; erosion rises in 77\% of trajectories & Iterative benchmark tasks; not SOLID-specific \\
\addlinespace
\citet{zhu2026smells} & LLM- and agent-driven projects & Code and architecture smells & Machine signature; Volume-Quality Inverse Law & Smells, not principles \\
\addlinespace
\citet{cynthia2026postmerge} & Agent bug-fix PRs in AIDev & SonarQube issues introduced by merged PRs & Agent differences vanish after churn normalisation; merge success does not reflect quality & Merged PRs only; bug fixes only \\
\addlinespace
\citet{sawada2026maintenance} & Agent PRs in AIDev & Post-merge maintenance & Humans do most maintenance; AI files mostly get feature extensions & Maintenance activity, not design structure \\
\addlinespace
\citet{xu2025productivity} & Copilot users & Productivity and review load & Core developers review 6.5\% more code, own output falls 19\% & Organisational study; no design metrics \\
\end{longtable}
\normalsize
\end{landscape}

\section{Critical Analysis and Motivation for This Work}
\label{sec:gaps-contributions}

The studies reviewed above show that coding agents are well evaluated on function-level and
issue-level tasks, and that structural problems in generated code are real. They also leave three
limits, which motivate the experiment of this thesis.

\paragraph{Real contributions at scale.} Existing evaluations either test functions and single
issues (HumanEval, SWE-bench) or assess design on small sets of tasks (ProjDevBench's 20 problems,
NITR's 21 probes) or on newly generated projects, and most of the studies use a single language. Few
look at what agents actually contribute to existing repositories, including contributions that
maintainers rejected or never decided on.

\paragraph{Unmeasured validity of the detectors.} Many studies assume that a metric or a tool
output captures a design principle. Language-model studies report accuracy on curated examples, and
tool-based studies report counts of issues without checking that each issue is a violation of a given
principle. How far a structural detector agrees with a reading of the same code is rarely measured.

\paragraph{Agent-level comparison.} Much of the literature varies the model and treats the agent
software as fixed. \citet{bensghaier2026harness} show that the agent software affects quality with
the model held fixed, but whether this reaches design quality has not been compared per principle
across several agents on a shared corpus. The studies built on the AIDev corpus look at merge rates,
issues in merged bug fixes and post-merge maintenance, not at SOLID adherence.

This thesis therefore runs an experiment that addresses these limits. It measures adherence to the
five SOLID principles in real pull requests written by three coding agents, taken from AIDev, in two
languages, Java and Python, and it measures how far each detector can be trusted. Chapter~\ref{ch:methodology}
describes the design of this experiment.
